\slajdid{09_2-s014}
\begin{frame}[label=09_2-s014]
\gusttekst\setlength{\abovedisplayskip}{3pt}\setlength{\belowdisplayskip}{3pt}\setlength{\abovedisplayshortskip}{2pt}\setlength{\belowdisplayshortskip}{2pt}\textbf{Drugi način pisanja - „svi naponi su pozitivni – označavanje suprotno od n-kanalnog, struja ide od sorsa ka drejnu“}
\[I_{Gp}=0\]
\[I_{Dp}=I_{Sp}\ge0\quad\textit{za}\quad V_{{\color{DEcrvena}SG}p}\ge V_{Tp}\quad(V_{Tp}>0)\]
\[V_{SDpsat}=\frac{(V_{SGp}-V_{Tp})}{1+\frac{(V_{SGp}-V_{Tp})}{L_pE_{Cp}}}=\frac{(V_{SGp}-V_{Tp})L_pE_{Cp}}{L_pE_{Cp}+(V_{SGp}-V_{Tp})}\quad(E_{Cp}>0)\]
Tranzistor vodi u triodnoj, omskoj, neaktivnoj oblasti, kada je
\[V_{SDp}\le V_{SDpsat}\]
i tada je njegova struja data izrazom
\[I_{Dp}=\frac{k_p}2\frac1{1+\frac{V_{SDp}}{L_pE_{Cp}}}\left(2V_{SDp}(V_{SGp}-V_{Tp})-V_{SDp}^2\right)\]
Slično kao da smo upotrebili oznake kao kod n- kanalnog ali pisali na primer $|V_{GSp}|$ itd…
\end{frame}
